
\section{Introduction to Natural Deduction}

\cite[][Chapter 16, pp. 110-2]{Mag25}
\cite[][Chapter 17, pp. 113-41]{Mag25}

\subsection{The very idea of natural deduction}

As much work as the partial truth tables help us save, we are still left with
the impression that the process should be more constrained and rule-based.
Ideally, we'd like to be able to evaluate the semantic properties of an
argument in TFL on the basis of formal practices alone. This is the basic
idea behind \emph{natural deduction}. We will use \emph{proofs} that allow us
to demonstrate the truth of some TFL sentences when we assume the truth of
other TFL sentences. In a sense these proofs leverage rules reflecting the
intuitions that we put to use last week in quickly working through our partial
truth tables to find counter-examples. But the con

But also, for potential tautologies and entailments with $n$ unique atomic
sentences, the size of a truth table is $2^n$ lines. A familar number to
computer scientists, but let's pause to consider how fast this number gets
very very big.

\twoitems{ 
  
  \begin{center}
    \begin{tabular}{c c}
      sentences & lines \\
      \hline
      1  & 2     \\
      2  & 4     \\
      3  & 8     \\
      4  & 16    \\
      5  & 32    \\
      6  & 64    \\
      7  & 128   \\
      8  & 256   \\
      9  & 512   \\
      10 & 1,024 \\
    \end{tabular}
  \end{center}
  
}{
  \begin{center}
    \begin{tabular}{c c}
      sentences & lines \\
      \hline
      11 & 2,048     \\
      12 & 4,096     \\
      13 & 8,192     \\
      14 & 16,384    \\
      15 & 32,768    \\
      16 & 655,36    \\
      17 & 131,072   \\
      18 & 262,144   \\
      19 & 524,288   \\
      20 & 1,1048,576\\
    \end{tabular}
  \end{center}
}

It is not easy to glean insights from such a document and questions about
\emph{why} an argument is correct become extremely obscure. We'd really like
to get both the result \emph{and} a story about its justification from the same
process. 1000+ line truth tables are just a bad machanism to provide that.

\subsection{Basic rules for TFL}

We begin with a basic definition of proofs (aka \emph{derivations}).

\begin{definition}[Proof]
  \label{def:formal-proof}
  A \emph{proof} in Truth Functional Logic is a sequence of sentences that:
  \begin{enumerate}
  \item Begin with zero or more \emph{premises}
  \item Introduce new sentence only when they can be made by applying
    \emph{rules} to lines already in the proof
  \item End with a \emph{conclusion} (the sentence that we want to prove)
  \end{enumerate}  
\end{definition}

We have two kinds of rules in a Fitch-style natural deduction system:
\begin{enumerate}
\item \textbf{Introduction Rules:} that allow us to work toward a conclusion
  containing the operator being introduced; and
\item \textbf{Elimination Rules:} that allow us to break down sentences that
  are already in the proof so that we can get at the parts.
\end{enumerate}  

Before we get into the rules themselves, we should pause for a minute and
discuss the styles and formats of the proofs themselves. First, we will
separate the premises from the rest of the proof with a horizontal line. Second,
we will number each line to make it easier to reference later in the proof. 
Finally, we will mark every line with a \emph{justification}. This will be the
rule applied along with the previous lines that it was applied to. We will mark
premises with their own ``rule'' PR, which just indicates that we take them to
be true by assumption.

\vspace{-0.5\baselineskip}
\begin{fitchproof}
\hypo {1} {P\eand Q} \by{PR}{}
\have {3} {P}        \ae{1}
\end{fitchproof}

In this simple proof, we show that we can conclude $P$ if we assume $P\eand Q$.

The simplest rule in the system is simply that we can repeat lines that we have
already seen.

\Needspace{5\baselineskip}
\begin{fitchrule}[Repitition]
  \label{def:repitition-rule}
  \leavevmode
  \vspace{-0.9\baselineskip}
  \begin{fitchproof}
    \have[m]      {1} \ephi
    \have[\vdots] {2} {\vdots}    
    \have[n]      {3} \ephi    \r{1}
  \end{fitchproof}
  \vspace{-0.9\baselineskip}
  where $m<n$
\end{fitchrule}

For each of the operators, we will want to introduce a pair of rules: one an
introduction rule and the other an elimination rule.

Conjunction has the most straightforward rules. When both conjuncts occur in a
proof, we can introduce a conjunction of them. When a conjunction occurs in a
proof, we can eliminate it for either of the conjuncts.

\Needspace{6\baselineskip}
\begin{fitchrule}[\eand Introduction]
  \label{def:and-introduction-rule}
  \leavevmode
  \vspace{-0.9\baselineskip}
  \begin{fitchproof}
    \have[m] {1} {\ephi}
    \have[n] {2} {\epsi}
    \have[ ] {3} {\ephi\eand\epsi} \ai{1,2}
  \end{fitchproof}
\end{fitchrule}

\Needspace{5\baselineskip}
\begin{fitchrule}[\eand Elimination]
  \label{def:and-introduction-rule}
  \leavevmode
  There are two rules for \eand elimination, one for each conjunct: \\
  \twoitems{
  \begin{fitchproof}
    \have[n] {1} {\ephi\eand\epsi}
    \have[ ] {2} {\ephi}              \ae{1}
  \end{fitchproof}
  }{
  \begin{fitchproof}
    \have[n] {1} {\ephi\eand\epsi}
    \have[ ] {2} {\epsi}              \ae{1}
  \end{fitchproof}
  }
\end{fitchrule}

With just the introduction and elimination rules for \eand we can prove many
valid arguments such as $A \eand B, C \eand D \therefore A \eand D$.

\begin{fitchproof}
  \hypo {1} {A \eand B} \by{PR}{}
  \hypo {2} {C \eand D} \by{PR}{}
  \have {3} {A}         \ae{1}
  \have {4} {D}         \ae{2}
  \have {5} {A \eand D} \ai{3,4}
\end{fitchproof}

The elimination rule is straightforward so we will start there.

\Needspace{6\baselineskip}
\begin{fitchrule}[\eif Elimination (Modus Ponens)]
  \label{def:if-elimination-rule}
  \leavevmode
  \vspace{-0.9\baselineskip}
  \begin{fitchproof}
    \have[m] {1} {\ephi\eif\epsi}
    \have[n] {2} {\ephi}
    \have[ ] {3} {\epsi}           \ie{1,2}
  \end{fitchproof}
\end{fitchrule}

So far, pretty simple. But for the introduction rule, we will need to make use
of assumptions that are not included in the premises. We will mark their
introduction with a new horizontal line and a justification of AS. We will
discuss assumptions in more detail after the introduction rule for conditionals.

\Needspace{6\baselineskip}
\begin{fitchrule}[\eif Introduction]
  \label{def:if-introduction-rule}
  \leavevmode
  \vspace{-0.9\baselineskip}
  \begin{fitchproof}
    \open
    \hypo[i] {1} {\ephi}           \by{AS}{}
    \have[j] {2} {\epsi}
    \close
    \have[ ] {3} {\ephi\eif\epsi}  \ii{1-2}
  \end{fitchproof}
\end{fitchrule}

The resulting sub-sequence, visibly set off from the rest of the proof using
indentation, is called a \emph{subproof}. When we use the \eif introduction rule
we collapse the subproof and treat its assumption as the antecedant of the new
conditional. At this point, the subproof is \emph{closed}.

Consider the following proof of $A\eif B, B \eif C \therefore A \eif C$.

\begin{fitchproof}
  \hypo {1} {A \eif B} \by{PR}{}
  \hypo {2} {B \eif C} \by{PR}{}
  \open
  \hypo {3} {A}        \by{AS}{}
  \have {4} {B}        \ie{1,3}
  \have {5} {C}        \ie{2,4}  
  \close
  \have {6} {A \eif C} \ii{3-5}
\end{fitchproof}

When the subproof is closed, the assumption that $A$ is true is set aside
(\emph{discharged}). We only needed it to establish that $A \eif C$, which is
true regardless of whether $A$ is. We should pause here and take a closer look
at this process.

Perhaps the most important point about subproofs is that their components are
meant to be used only by the subproofs themselves, not by the larger proof(s)
containing them. An open subproof must be closed; we cannot reach inside closed
subproofs.

\begin{definition}[Permissible Line Citations]
  \label{def:permissible-line-citations}
  When we want to justify a rule in a formal proof, we can only cite lines:
  \vspace{-0.9\baselineskip}
  \begin{enumerate}
  \item that occur before the current line
  \item that do not occur in a closed subproof
  \end{enumerate}  
\end{definition}

Similarly, we want to be careful about only citing closed subproofs when we
apply rules like \ref{def:if-introduction-rule}. We can only cite closed
subproofs and they cannot be nested too deeply.

\begin{definition}[Permissible Subproof Citations]
  \label{def:permissible-subproof-citations}
  When we want to justify a rule in a formal proof, we can only cite subproofs:
  \vspace{-0.9\baselineskip}
  \begin{enumerate}
  \item that occur entirely before the current line
  \item that do not occur inside another closed subproof
  \item whose last line does not occur inside a nested subproof
  \end{enumerate}  
\end{definition}

The third condition is really just a special case of the second.

With this basic understanding of subproofs and definitions for \eif
introduction and \eif elimination, the rules for biconditionals are trivial. 

\Needspace{5\baselineskip}
\begin{fitchrule}[\eiff Elimination]
  \label{def:iff-elimination-rule}
  \leavevmode
  %\vspace{-0.9\baselineskip}
  There are two rules for \eiff elimination, one for each sentence: \\
  \twoitems{
    \begin{fitchproof}
      \have[m] {1} {\ephi\eiff\epsi}
      \have[n] {2} {\ephi}
      \have[ ] {3} {\epsi}           \by{\eiff E}{1,2}
    \end{fitchproof}
  }{
    \begin{fitchproof}
      \have[m] {1} {\ephi\eiff\epsi}
      \have[n] {2} {\epsi}
      \have[ ] {3} {\ephi}           \by{\eiff E}{1,2}
    \end{fitchproof}
  }
\end{fitchrule}

\Needspace{6\baselineskip}
\begin{fitchrule}[\eiff Introduction]
  \label{def:iff-introduction-rule}
  \leavevmode
  \vspace{-0.9\baselineskip}
  \begin{fitchproof}
    \open
    \hypo[i] {1} {\ephi}           \by{AS}{}
    \have[j] {2} {\epsi}
    \close
    \open
    \hypo[k] {3} {\epsi}           \by{AS}{}
    \have[l] {4} {\ephi}
    \close
    \have[ ] {5} {\ephi\eiff\epsi} \by{\eiff I}{1-2,3-4}
  \end{fitchproof}
\end{fitchrule}

The last few operators are more challenging to defnine. The first of these,
disjunction (\eor), because its elimination requires us to check more than one
possibility.

Introduction, however, is quite straightforward.

\Needspace{5\baselineskip}
\begin{fitchrule}[\eor Introduction]
  \label{def:eor-introduction-rule}
  \leavevmode
  %\vspace{-0.9\baselineskip}
  There are two rules for \eor introduction, one for each sentence: \\
  \twoitems{
    \begin{fitchproof}
      \have[n] {1} {\ephi}
      \have[ ] {2} {\ephi\eor\epsi}  \oi{1}
    \end{fitchproof}
  }{
    \begin{fitchproof}
      \have[n] {1} {\epsi}
      \have[ ] {2} {\ephi\eor\epsi}  \oi{1}
    \end{fitchproof}
  }
\end{fitchrule}

We can handle the ambiguity involved in disjunction elimination with multiple
subproofs, much as we do when intoducing biconditionals.

\Needspace{8\baselineskip}
\begin{fitchrule}[\eor Elimination]
  \label{def:eor-elimination-rule}
  \leavevmode
  \vspace{-0.9\baselineskip}
  \begin{fitchproof}
    \have[j] {1} {\ephi\eor\epsi}
    \open
    \hypo[i] {2} {\ephi}  \by{AS}{}
    \have[j] {3} {\echi}
    \close
    \open
    \hypo[k] {4} {\epsi}  \by{AS}{}
    \have[l] {5} {\echi}
    \close
    \have[ ] {6} {\echi}  \oe{1, 2-3,4-5}
  \end{fitchproof}
\end{fitchrule}

Now that we have almost a complete set of rules for the basic operators of TFL,
we can prove a variety of valid arguments, such as this one demonstrating the
distributivity of \eand and \eor.

\begin{fitchproof}
  \hypo {1} {A \eand (B \eor C)}          \by{PR}{}
  \have {2} {A}                           \ae{1}
  \have {3} {B \eor C}                    \ae{1}
  \open
  \hypo {4} {B}                           \by{AS}{}
  \have {5} {A \eand B}                   \ai{2,4}
  \have {6} {(A \eand B)\eor(A \eand C)}  \oi{5}
    \close
  \open
  \hypo {7} {C}                           \by{AS}{}
  \have {8} {A \eand C}                   \ai{2,7}
  \have {9} {(A \eand B)\eor(A \eand C)}  \oi{8}
  \close
  \have {10} {(A \eand B)\eor(A \eand C)} \oe{3,4-6,7-9}
\end{fitchproof}

An attentive reader will probably have noticed that we have not yet mentioned
the negation operator (\enot). The nature of negation in TFL is wrapped up in
two of its core assumptions: \emph{non-contradiction} and \emph{the excluded
middle}. The first says that \ephi and \enot\ephi cannot both be true:
\enot(\ephi\eand\enot\ephi). The second says that one of either \ephi or
\enot\ephi must be true: \ephi\eor\enot\ephi. To accomomdate these assumptions,
especially the first, we introduce an additional symbol into our object
language for TFL (\ebot) which we will use when we encounter a contradiction.


\Needspace{4\baselineskip}
\begin{fitchrule}[\enot Elimination]
  \label{def:not-elimination-rule}
  \leavevmode
  \vspace{-0.9\baselineskip}
  \begin{fitchproof}
    \have[m] {1} {\enot\ephi}
    \have[n] {2} {\ephi}
    \have[ ] {3} {\ebot}      \ne{1,2}
  \end{fitchproof}
\end{fitchrule}

There are several ways to think about \ebot but we will just assume that it is
a sentence that is always false, regardless of situation. This makes it easy
to use in other rules and does not obscure its use to indicate a contradiction.

But how do we introduce a negation? We assume the non-negated sentence and use
a subproof to show that this assumption leads to a contradiction. The
introduction rule for negation references subproofs of this form in its
justification.

\Needspace{4\baselineskip}
\begin{fitchrule}[\enot Introduction]
  \label{def:not-introduction-rule}
  \leavevmode
  \vspace{-0.9\baselineskip}
  \begin{fitchproof}
    \open
    \hypo[m] {1} {\ephi}      \by{AS}{}
    \have[n] {2} {\ebot}
    \close
    \have[ ] {3} {\enot\ephi} \ni{1-2}
  \end{fitchproof}
\end{fitchrule}

Of course, the symmetry of negation lets us work in the opposite direction,
too. We will call this method: \emph{indirect proof}, since it does not involve
the introduction or elimation of any operators outside of the subproof.

\Needspace{4\baselineskip}
\begin{fitchrule}[Indirect Proof]
  \label{def:indirect-proof-rule}
  \leavevmode
  \vspace{-0.9\baselineskip}
  \begin{fitchproof}
    \open
    \hypo[m] {1} {\enot\ephi} \by{AS}{}
    \have[n] {2} {\ebot}
    \close
    \have[ ] {3} {\ephi}      \by{IP}{1-2}
  \end{fitchproof}
\end{fitchrule}

One additional rule can be thought of as \ebot elimination. A long-standing
assumption in classical logic is that proving a contradiction allows us to
conclude anything at all.

\Needspace{3\baselineskip}
\begin{fitchrule}[Explosion]
  \label{def:explosion-rule}
  \leavevmode
  \vspace{-0.9\baselineskip}
  \begin{fitchproof}
    \have[m] {1} {\ebot}
    \have[ ] {2} {\ephi} \by{X}{1}
  \end{fitchproof}
\end{fitchrule}
